\documentclass{article}
\usepackage[UTF8]{ctex}
\usepackage{amsmath,amssymb,amsfonts}
\usepackage{enumitem}
\usepackage{geometry}
\geometry{a4paper, margin=1in}
% === 错题本排版优化（不参与题目标识） ===
\setlength{\parindent}{0pt}
\setlength{\parskip}{0.35em plus 0.10em minus 0.08em}
\setlength{\emergencystretch}{3em}
\setlist{itemsep=0.25em, topsep=0.40em, parsep=0pt}
\allowdisplaybreaks[2]
% === 错题本排版优化结束 ===
\begin{document}
\section*{西南交通大学2024年数学分析真题}
\section*{170．西南交通大学2024年数学分析试卷（回忆版）}

\subsection{题目1}
1．已知 \(\lim _{n \rightarrow \infty} a_{n}=a\) ，证明： \(\lim _{n \rightarrow \infty} \frac{\left[n a_{n}\right]}{n}=a\) ，其中［•］表示取整。

\subsection{题目2}
2．设 \(f(x)\) 在 \((a, b)\) 上连续，且 \(f(a+0), f(b-0)\) 存在，证明：\(f(x)\) 在 \((a, b)\) 上一致连续．

\subsection{题目3}
3．设 \(f(x)\) 在 \([a, b]\) 上连续，在 \((a, b)\) 上可导，证明：存在 \(\xi \in(a, b)\) ，使得 \(f^{\prime}(\xi)+\lambda f(\xi)=0\) ．

\subsection{题目4}
5．设 \(f(x)\) 在 \([a, b]\) 上有界，判断 \(f,|f|, f^{2}\) 之间的可积性关系．

\subsection{题目5}
6．讨论级数 \(\sum_{n=1}^{\infty} \frac{\sin n x}{n^{p}}\) 的敛散性。

\subsection{题目6}
7．证明 \(f(x, y)=\sqrt{|x y|}\) 在点 \((0,0)\) 处连续，并求 \(f_{x}(0,0), f_{y}(0,0)\) ，判断 \(f(x, y)\) 在点 \((0,0)\) 的可微性。

\subsection{题目7}
8．设 \(f(x)\) 在 \((0,+\infty)\) 上连续，当 \(x \rightarrow 0^{+}\)时，\(f(x) \rightarrow+\infty\) ，且 \(\int_{0}^{1} f(x) \mathrm{d} x\) 收敛，证明： \(\lim _{x \rightarrow 0^{+}} x f(x)=0\) ．

\subsection{题目8}
9．设 \(u=f(x, y)\) ，其中 \(x=r \cos \theta, y=r \sin \theta\) ，且 \(\frac{\partial u}{\partial x}=\frac{\partial v}{\partial y}, \frac{\partial u}{\partial y}=-\frac{\partial v}{\partial x}\) ，证明：

\[
\frac{\partial^{2} u}{\partial x^{2}}=\frac{1}{r^{2}} \frac{\partial^{2} v}{\partial y^{2}}, \frac{\partial^{2} u}{\partial y^{2}}=-\frac{1}{r^{2}} \frac{\partial^{2} v}{\partial x^{2}}
\]

\subsection{题目9}
10．求 \(y=m x^{2}, y=n x^{2}(0<m<m), y=\alpha x, y=\beta x(0<\alpha<\beta)\) 所围区域的面积．

\subsection{题目10}
11．设 \(\Omega\) 是单连通区域，\(L\) 为其边界，在 \(\Omega\) 内部或外部取一定点，设其为原点，\(r\) 为 \(L\) 上的点到原点的向量，\(n\) 为该点的单位切向量，求 \(\oint_{L} \frac{\cos (n, r)}{|r|} \mathrm{d} s\) ．

\subsection{题目11}
12．设 \(f(x)\) 在 \([a, b]\) 上无界，证明：存在 \(x_{0} \in(a, b)\) ，使得对任意的 \(\delta>0\) ，满足 \(f(x)\) 在\(\left(x_{0}-\delta, x_{0}+\delta\right) \cap[a, b]\) 上无界．
\end{document}
